\section{第二问：变物性热质耦合与全过程输出}
\label{sec:q2}

\subsection{状态反馈与全程求解}
第二问从共同初态开始，全程使用附录3物性\cite{problem}，不接续第一问
末态，也不在$1800\unit{s}$处切换公式。有效热容量
$S(C)=\rho_{\rm eff}(C)c_p(C)$和导热率随局部含水率改变，扩散率同时
依赖局部绝对温度$T_K=T+273.15$和$C$。升温改变扩散，失水又通过热容量
及导热率反馈温度，因此需联立推进两个场。

共用方程中的$k$、$D$均留在散度内。特别地，
$S^{-1}\nabla\cdot(k\nabla T)$一般不等于
$\nabla\cdot[(k/S)\nabla T]$，不能把变热扩散率直接当作热通量系数。
水分方程利用可约去的参考干密度解释；经验$\rho_{\rm eff}(C)$在这里
用于有效热容量，并不同时充当固定体积下满足真实组成守恒的湿体密度。

前$4\unit{h}$采用附件1分段线性环境，之后按统一假设保持末小时61点
算术均值：$T_a=49.9989344262\celsius$、$b=0.04998754098$。
第二问与第三问采用同一条温湿联合轨迹，积分至第三问执行端点
$206906.76\unit{s}$，即$57.4741\unit{h}$。因此第二问的完整结果
覆盖所选模型下的全过程，题设表3、表4对应的表\ref{tab:q2-temperature}、\ref{tab:q2-moisture}仅展示前$3\unit{h}$。
环境平台是超出观测范围的假设，不能将此完整时间覆盖称为数天实测验证。

全程轨迹采用80阶Chebyshev配点、原函数形式水分通量和Radau隐式
时间积分，表面温湿由Robin条件联合确定；相对容差$2\times10^{-13}$，
最大步长$30\unit{s}$，前$4\unit{h}$限制为$20\unit{s}$。
\texttt{result2.xlsx}包含两张表，每表从1秒起保存206907个时刻，即
整数秒至206906秒并追加206906.76秒；21个规定半径合计8,690,094个
温湿数值。逐秒评价联合连续解，不将不同模型的温度、水分拼接。

\subsection{规定表格及空间规律}
\begin{table}[htbp]
\centering\small
\caption{3小时内药材中截面温度（$^\circ$C）}
\label{tab:q2-temperature}
\input{tables/q2_temperature.tex}
\end{table}

\begin{table}[htbp]
\centering\small
\caption{3小时内药材中截面干基含水率（kg/kg）}
\label{tab:q2-moisture}
\input{tables/q2_moisture.tex}
\end{table}

至$3\unit{h}$，轴心、表面温度分别为49.8495\celsius、49.9664\celsius，
相差$0.1169\unit{K}$；含水率分别为1.7662、1.0081，相差
$0.7581\unit{kg/kg}$。温度接近空间均匀时，水分仍有明显梯度，
不能用均匀含水率的集中参数近似替代径向场。
按$r\,\mathrm dr$加权，平均温度为49.9043\celsius，平均含水率为
$1.3825\unit{kg/kg}$，相对初始含水率下降45.7833\%。该比例是固定参考
干密度下的有效水分减少比例，不是实测总质量损失率。

\begin{figure}[htbp]
\centering
\includegraphics[width=0.94\textwidth]{figures/q2_profiles.pdf}
\caption{第二问前3小时的中截面径向温湿分布。}
\label{fig:q2-profiles}
\end{figure}

图\ref{fig:q2-profiles}表明，初期失水集中于外层，随后内部含水率逐步
降低。第二问$0.5\unit{h}$轴心温度32.1892\celsius，低于第一问同一
时刻；这是从初态使用不同物性、尤其初始有效热容量更大的正常结果，
不是问间接口错误。第一问和第二问是不同参数情形，不能共用一段预热轨迹。

\subsection{升温促进与失水抑制的竞争}
将附录3扩散系数相对初态取对数，可精确分解为
\begin{equation}
\begin{aligned}
\ln\frac{D(t,r)}{D_0}&=H(t,r)+M(t,r),\\
H&=3850\left(\frac1{301.15}-\frac1{T_K}\right),\qquad
M=0.45\left(\frac1{2.55}-\frac1C\right).
\end{aligned}
\label{eq:q2-log-mechanism}
\end{equation}
其中$D_0=5.64168037\times10^{-9}\unit{m^2/s}$。
$H$度量沿当前轨迹的累计温度贡献，$M$度量累计含水率贡献；这是代数
分解，不能当作独立改变某一因素的因果贡献百分比。

$3\unit{h}$时，轴心与表面的$D/D_0$分别为2.1957与1.8207，说明相对
初态累计升温促进仍占优势；但最后半小时，两处扩散率分别下降约1.06\%
和3.16\%。因此“扩散系数全程单调增加”不成立。表面还曾在早期每秒
采样的第144秒达到$D/D_0=0.983282$的局部最低值。温度变化和失水抑制
的竞争随阶段变化，前3小时的累计主导关系不能直接用于判定末期干燥瓶颈。

在同一320区间有限体积网格上单独扰动交换系数，$h_m$降低、提高20\%
时，$3\unit{h}$表面含水率分别为1.1775、0.8737；相应轴心值为1.8581、
1.6921。对这一输出和扰动区间，水分交换系数的影响明显大于数值误差。
这是人为设定的确定性情景，不是参数置信区间，也不能用它替代空气量
与材料平衡含水率的基准识别\cite{dasilva2014,adrover2020}。

\subsection{独立验证与适用范围}
前3小时已有节点有限体积逐级加密至2560区间及Richardson外推，
另用独立空间离散与Radau积分核对。全程轨迹与该外推轨迹在前3小时
共同采样点的最大温差约$1.07\times10^{-9}\unit{K}$、含水率差约
$8.41\times10^{-8}\unit{kg/kg}$；前3小时温湿表及规定60个论文表值四位一致。
有限体积外推是这一前3小时交叉验证的来源，不是完整工作簿的生成方法。

长时段另与已完成的160阶BDF轨迹在全部60秒采样点及同一执行端点比较，
最大温湿差约$3.27\times10^{-9}\unit{K}$与$9.33\times10^{-9}\unit{kg/kg}$，
四位输出一致。当前869万余格工作簿与同源未舍入轨迹的舍入核对全部通过，
该检查支持导出一致性；高阶独立参考覆盖的是所列采样点，不能把导出
一致性误写成全部逐秒状态均另作了独立求解。

原有限体积有效水分账本的最大残差约$6.51\times10^{-12}\unit{kg/kg}$。
热账本核对的是有效热容量积分与边界热流，即比较
$\int\!\int S(C)T_t\,\mathrm dV\mathrm dt$与累计边界热量，
相对残差约$3.11\times10^{-11}$，并非$\int ST\,\mathrm dV$的末初差，
更不构成潜热、携焓及真实总能量已经闭合的证明。
同物理$80\times256$二维配对在前3小时已检查采样点的中截面最大温湿差
约为$1.692\times10^{-3}\unit{K}$、$2.247\times10^{-5}\unit{kg/kg}$。
这支持中截面的工程近似，但不能保证降维后所有四位显示相同；长期全域
达标须结合第三问终点证据，不能直接沿用前3小时的几何误差。
\FloatBarrier
